\input{../tex-inputs/latex-leseansicht-vorspann}

\section[Arthur und Olga Schnitzler an Gustav Schwarzkopf, 12. 8. 1911]{L04388 Arthur und Olga Schnitzler an Gustav Schwarzkopf, 12. 8. 1911}
\nopagebreak\mylabel{L04388v}
\rehead{ }\normalsize\beginnumbering\briefempfaengerindex{Schwarzkopf, Gustav@\textsc{Schwarzkopf, Gustav}!zzzSchnitzler, Olga@\emph{von Olga Schnitzler}!1911-08-121@{12. 8. 1911}|(be}\briefempfaengerindex{Schwarzkopf, Gustav@\textsc{Schwarzkopf, Gustav}!zzzSchnitzler, Arthur@\emph{von Arthur Schnitzler}!1911-08-121@{12. 8. 1911}|(be}
\toendnotes[C]{\smallbreak\pagebreak[2]}
\correspDesc{Versand  durch Arthur Schnitzler, Olga Schnitzler am 12. 8. 1911 in Semmering
\newline{}Übermittlung  am 12. 8. 1911 in Semmering
\newline{}Erhalt  durch Gustav Schwarzkopf im Zeitraum [13. 8. 1911 – 17. 8. 1911?] in Wien}\toendnotes[C]{\smallbreak}
\Standort{DLA, A:Schnitzler, HS.1985.1.1897.}
\physDesc{Bildpostkarte, 620 Zeichen
\newline{}Handschrift Arthur Schnitzler: Bleistift, deutsche Kurrent
\newline{}Handschrift Olga Schnitzler: Bleistift, lateinische Kurrent
\newline{}Versand: Stempel: »\nobreak{}\oindex{Semmering@\textbf{Semmering}|pwk}Semmering, 12 VIII 11, 3\nobreak{}«.  }\pstart{}{\pb}Herrn \textsc{Gustav
                     Schwarzkopf}\pend{}\pstart{}\textsc{Wien I}\oindex{I., Innere Stadt@\textbf{I., Innere Stadt}|pw}\pend{}\pstart{}\textsc{Tiefer Graben 23}\oindex{Wien@\textbf{Wien}!I., Innere Stadt@\textbf{I., Innere Stadt}!Tiefer Graben 23@\textbf{Tiefer Graben 23}|pw}.\pend{}{\bigskip}
\pstart
           \centering{}{\pb}\textcolor{gray}{\textbf{Semmering. Südbahnhotel}}\oindex{Südbahnhotel [Semmering]@\textbf{Südbahnhotel [Semmering]}|pw}\pend
           \vspace{1em}
\pstart
           \noindent{}{\pb}lieber Guſtav, wir fühlen uns jetzt hier recht wohl, bleiben 2–3
               Wochen; und es wäre ſehr hübſch, we{\geminationn}{ }ſie einmal oder
               öfters oder auf länger heraufkämen. Von Richard\pwindex{Beer-Hofmann, Richard 11.\,7.\,1866 Wien – 26.\,9.\,1945 New York City@\textsc{Beer-Hofmann, Richard} (11.\,7.\,1866 Wien – 26.\,9.\,1945 New York City)|pw} weiſs ich nur, daſs er vor circa 14 Tagen in \textsc{Ischl\oindex{Bad Ischl@\textbf{Bad Ischl}|pw}} mit \textsc{S. Fischer}’s\pwindex{Fischer, Samuel 24.\,12.\,1859 Liptovský Mikuláš – 15.\,10.\,1934 Berlin@\textsc{Fischer, Samuel} (24.\,12.\,1859 Liptovský Mikuláš – 15.\,10.\,1934 Berlin)|pw}\pwindex{Fischer, Hedwig 8.\,9.\,1871 Szczecin – 11.\,4.\,1952 Königstein im Taunus@\textsc{Fischer, Hedwig} (8.\,9.\,1871 Szczecin – 11.\,4.\,1952 Königstein im Taunus)|pw} u \textsc{Salten}’s\pwindex{Salten, Felix 6.\,9.\,1869 Budapest – 8.\,10.\,1945 Zürich@\textsc{Salten, Felix} (6.\,9.\,1869 Budapest – 8.\,10.\,1945 Zürich)|pw}\pwindex{Salten, Ottilie 7.\,3.\,1868 Prag – 22.\,6.\,1942 Zürich@\textsc{Salten, Ottilie} (7.\,3.\,1868 Prag – 22.\,6.\,1942 Zürich)|pw} Rendezvous hatte. Directe Nachricht hab
               ich nicht. Jedenfalls iſt er noch in Auſſee\oindex{Bad Aussee@\textbf{Bad Aussee}|pw}. Von
                  Beka{\geminationn}ten iſt nur \textsc{Hugo Salus\pwindex{Salus, Hugo 3.\,8.\,1866 Česká Lípa – 4.\,2.\,1929 Prag@\textsc{Salus, Hugo} (3.\,8.\,1866 Česká Lípa – 4.\,2.\,1929 Prag)|pw}} hier. Ich arbeite nur wenig. Laſſen Sie bald von{ }ſich hören oder laſſen noch{ }ſich{ }ſehen\pend
           \pstart Wir grüßen herzlichſt. \spacefill\mbox{Ihr A. }\pend{}\selectlanguage{ngerman}\vspace{1em}
\pstart
           \noindent{}{[}hs. O. Schnitzler:{]} Bitte kommen Sie, es ist wunderschön hier und Lili\pwindex{Cappellini, Lili 13.\,9.\,1909 Wien – 26.\,7.\,1928 Venedig@\textsc{Cappellini, Lili} (13.\,9.\,1909 Wien – 26.\,7.\,1928 Venedig)|pw} wird Sie sehr amüsieren!\pend
           
\pstart
           Mit herzl. Grüssen{\\[\baselineskip]}Ihre{\\[\baselineskip]}\spacefill\mbox{O. S.}\pend
           \leftskip=0em{}
\pstart
           12/8 911\pend
           \selectlanguage{ngerman}\endnumbering\briefempfaengerindex{Schwarzkopf, Gustav@\textsc{Schwarzkopf, Gustav}!zzzSchnitzler, Olga@\emph{von Olga Schnitzler}!1911-08-121@{12. 8. 1911}|)be}\briefempfaengerindex{Schwarzkopf, Gustav@\textsc{Schwarzkopf, Gustav}!zzzSchnitzler, Arthur@\emph{von Arthur Schnitzler}!1911-08-121@{12. 8. 1911}|)be}\mylabel{L04388h}
\begin{anhang}
\end{anhang}\newcommand{\dateiname}{L04388}\newcommand{\titel}{Arthur und Olga Schnitzler an Gustav Schwarzkopf, 12. 8. 1911}\newcommand{\editorInnen}{Herausgegeben von Jahnke, SelmaMüller, Martin Anton}\input{../tex-inputs/latex-leseansicht-abspann}
